\documentclass[12pt,a4paper]{article}

\usepackage[utf8]{inputenc}
\usepackage[T1]{fontenc}
\usepackage[spanish]{babel}
\usepackage{geometry}
\usepackage{hyperref}
\usepackage{listings}
\usepackage{xcolor}
\usepackage{graphicx}
\usepackage{tabularx}
\usepackage{enumitem}
\usepackage{fancyhdr}

\geometry{margin=2.5cm}

\definecolor{codebg}{rgb}{0.95,0.95,0.95}
\definecolor{codeblue}{rgb}{0.2,0.2,0.6}
\definecolor{codegreen}{rgb}{0.1,0.5,0.1}

\lstset{
    backgroundcolor=\color{codebg},
    basicstyle=\ttfamily\small,
    keywordstyle=\color{codeblue},
    commentstyle=\color{codegreen},
    frame=single,
    breaklines=true,
    tabsize=2,
    captionpos=b
}

\hypersetup{
    colorlinks=true,
    linkcolor=blue,
    urlcolor=blue
}

\pagestyle{fancy}
\lhead{Vetrinaria}
\rhead{Sprint 0 y 1}
\cfoot{\thepage}

\title{\vspace{-2cm}\textbf{Vetrinaria}\\[0.3cm]
\large Sistema de Gestión para Clínicas Veterinarias\\[0.5cm]
\textbf{Documentación Técnica — Sprint 0 y 1}}
\author{}
\date{Mayo 2026}

\begin{document}
\maketitle
\thispagestyle{empty}
\newpage

\tableofcontents
\newpage

\section{Resumen Ejecutivo}

Este documento detalla el trabajo realizado durante los Sprint 0 y 1 del proyecto
\textbf{Vetrinaria}, un sistema integral para la gestión de clínicas veterinarias.
El objetivo de estos sprints fue establecer la arquitectura base del proyecto y el
sistema de autenticación.

\begin{itemize}[leftmargin=*]
    \item \textbf{Sprint 0:} Configuración del monorepo, base de datos, Docker,
          componentes UI y diseño responsivo.
    \item \textbf{Sprint 1:} Sistema de autenticación con NextAuth v5, páginas de
          login/registro y middleware de protección por roles.
\end{itemize}

\section{Sprint 0 — Setup del Proyecto}

\subsection{Arquitectura del Monorepo}

El proyecto utiliza \textbf{Turborepo} con \textbf{pnpm} como gestor de paquetes,
organizado en la siguiente estructura:

\begin{lstlisting}[language=bash, caption=Estructura del monorepo]
vetrinaria/
+-- apps/
|   +-- web/                    # Next.js 15 App Router
|       +-- src/
|           +-- app/            # Rutas y layouts
|           +-- components/     # Componentes React
|           +-- lib/            # Utilidades y config
+-- packages/
|   +-- db/                     # Drizzle ORM + esquemas
|   +-- shared/                 # Tipos y constantes
|   +-- config/                 # ESLint, tsconfig base
+-- docker-compose.yml          # Servicios PostgreSQL, MinIO, Redis
+-- turbo.json                  # Configuración de Turborepo
\end{lstlisting}

Las dependencias se gestionan con \texttt{pnpm-workspace.yaml}, que define los
workspaces del monorepo. Turborepo permite ejecutar tareas (\texttt{build},
\texttt{lint}, \texttt{typecheck}) en paralelo con caching inteligente.

\subsection{Configuración de Next.js 15 y TypeScript}

Se configuró \textbf{Next.js 15} con el \textbf{App Router} y \textbf{TypeScript 5.8},
utilizando configuraciones base compartidas desde \texttt{packages/config}:

\begin{itemize}[leftmargin=*]
    \item \texttt{tsconfig.base.json}: Configuración base con ES2022, strict mode,
          y bundler module resolution.
    \item \texttt{tsconfig.nextjs.json}: Extiende la base con jsx:preserve y
          el plugin de Next.js.
\end{itemize}

El build de Next.js 15 compila exitosamente en aproximadamente 30 segundos
en el entorno de desarrollo.

\subsection{Base de Datos — PostgreSQL y Docker Compose}

Se implementaron tres servicios en Docker Compose:

\begin{itemize}[leftmargin=*]
    \item \textbf{PostgreSQL 16:} Puerto 5433 (evitando conflicto con instalación
          del sistema en 5432), con usuario \texttt{vet} y base de datos
          \texttt{vetrinaria}.
    \item \textbf{MinIO:} Almacenamiento de objetos S3-compatible para archivos
          (imágenes, radiografías, facturas).
    \item \textbf{Redis:} Caché y manejo de sesiones.
\end{itemize}

Todos los contenedores se inician con \texttt{docker compose up -d} y comparten
una red Docker interna.

\subsection{Esquema de Base de Datos — Drizzle ORM}

Se definieron 7 tablas usando \textbf{Drizzle ORM} con PostgreSQL:

\begin{table}[h]
\centering
\caption{Entidades del esquema de base de datos}
\begin{tabularx}{\textwidth}{|l|X|l|}
\hline
\textbf{Tabla} & \textbf{Propósito} & \textbf{Campos clave} \\
\hline
\texttt{users} & Usuarios del sistema & email, password\_hash, role \\
\texttt{clients} & Clientes (dueños de mascotas) & name, email, phone \\
\texttt{patients} & Pacientes (mascotas) & name, species, breed, client\_id \\
\texttt{appointments} & Citas veterinarias & date, status, patient\_id, user\_id \\
\texttt{medical\_records} & Historial clínico & diagnosis, treatment, patient\_id \\
\texttt{products} & Productos e inventario & name, price, stock \\
\texttt{suppliers} & Proveedores & name, contact, phone \\
\hline
\end{tabularx}
\end{table}

Las relaciones clave incluyen:

\begin{itemize}[leftmargin=*]
    \item \texttt{patients} → \texttt{clients} (FK: \texttt{client\_id})
    \item \texttt{appointments} → \texttt{patients}, \texttt{clients}, \texttt{users}
    \item \texttt{medical\_records} → \texttt{patients}, \texttt{appointments}, \texttt{users}
\end{itemize}

\subsection{TailwindCSS v4 y shadcn/ui}

Se implementó \textbf{TailwindCSS v4} con \textbf{shadcn/ui} para los componentes
de interfaz. Los componentes base creados son:

\begin{itemize}[leftmargin=*]
    \item \textbf{Button:} Variantes (default, destructive, outline, secondary,
          ghost, link) y tamaños (default, sm, lg, icon).
    \item \textbf{Card:} Contenedor con header, content y footer.
    \item \textbf{Input:} Campo de texto con validación.
    \item \textbf{Label:} Etiqueta accesible para formularios.
    \item \textbf{Avatar:} Imagen circular con fallback de iniciales.
    \item \textbf{Separator:} Línea divisoria horizontal.
    \item \textbf{DropdownMenu:} Menú contextual desplegable.
\end{itemize}

Todos los componentes siguen la convención \texttt{cn()} para fusión de clases
Tailwind, usando \texttt{clsx} y \texttt{tailwind-merge}.

\subsection{Layout Responsivo}

Se implementó un layout de dos columnas:

\begin{itemize}[leftmargin=*]
    \item \textbf{Sidebar:} Colapsable en móviles, con navegación principal.
          Utiliza estado global para el toggle y se cierra automáticamente al
          navegar en dispositivos móviles.
    \item \textbf{Navbar:} Barra superior con notificaciones, cambio de tema
          (claro/oscuro) y menú de usuario.
    \item \textbf{ThemeProvider:} Implementado con \texttt{next-themes}, soporta
          tema claro, oscuro y seguimiento del sistema.
\end{itemize}

El layout usa \textbf{route groups} de Next.js para separar la interfaz del
dashboard (con sidebar/navbar) de las páginas de autenticación (sin ellos).

\subsection{CI/CD — GitHub Actions}

Se configuró un pipeline de CI en \texttt{.github/workflows/ci.yml} que ejecuta
en cada push a \texttt{main} y cada Pull Request:

\begin{enumerate}[leftmargin=*]
    \item \textbf{Lint:} Verificación de código con ESLint.
    \item \textbf{Typecheck:} Validación de tipos con TypeScript.
    \item \textbf{Build:} Compilación de producción.
\end{enumerate}

\section{Sprint 1 — Autenticación y Roles}

\subsection{NextAuth v5}

Se implementó \textbf{NextAuth v5} (beta 31) con el proveedor de credenciales
(email y contraseña) y estrategia de sesión JWT:

\begin{lstlisting}[language=TypeScript, caption=Configuración de NextAuth]
export const { handlers, signIn, signOut, auth } = NextAuth({
  providers: [
    Credentials({
      async authorize(credentials) {
        // Validar contra BD con bcrypt
        const user = await db.select().from(schema.users)
          .where(eq(schema.users.email, email))
          .then((rows) => rows[0]);

        const isValid = await compare(password, user.passwordHash);
        if (!isValid) return null;

        return { id: String(user.id), email: user.email,
                 name: ..., role: user.role };
      },
    }),
  ],
  session: { strategy: 'jwt' },
});
\end{lstlisting}

Los tipos de TypeScript se extendieron para incluir \texttt{role} e \texttt{id}
en los objetos \texttt{User}, \texttt{Session} y \texttt{JWT}.

\subsection{Páginas de Autenticación}

Se crearon dos páginas en la ruta \texttt{/auth}:

\subsubsection*{Login (\texttt{/auth/login})}
Formulario con campos de email y contraseña. Utiliza el método
\texttt{signIn('credentials')} del lado del cliente. Muestra errores de
credenciales inválidas y redirige al dashboard tras autenticación exitosa.

\subsubsection*{Registro (\texttt{/auth/register})}
Formulario con nombre, apellido, email y contraseña. Envía los datos a un
endpoint API (\texttt{POST /api/auth/register}) que hashea la contraseña con
bcrypt (12 rondas) y crea el usuario en la base de datos con rol
\texttt{receptionist} por defecto.

\subsection{Middleware de Protección}

Se implementó un middleware que protege las rutas según el estado de
autenticación y el rol del usuario:

\begin{itemize}[leftmargin=*]
    \item \textbf{Rutas públicas:} \texttt{/auth/login}, \texttt{/auth/register}
    \item \textbf{Ruta raíz:} Redirige a login si no hay sesión, al dashboard
          si la hay.
    \item \textbf{Protección por rol:} El middleware clasifica las rutas
          accesibles para cada rol (admin y veterinarian tienen acceso completo;
          technician y receptionist tienen acceso restringido).
\end{itemize}

\subsection{Integración con el Layout}

La sesión de usuario se integra en la interfaz a través de:

\begin{itemize}[leftmargin=*]
    \item \textbf{SessionProvider:} Envuelve toda la aplicación en el root layout.
    \item \textbf{Navbar:} Muestra el nombre del usuario autenticado, sus
          iniciales en el avatar, y permite cerrar sesión.
    \item \textbf{Dashboard:} Verifica la sesión en el servidor y redirige
          a login si no está autenticado.
\end{itemize}

\subsection{Manejo de Errores}

El endpoint de registro valida:

\begin{itemize}[leftmargin=*]
    \item Todos los campos son obligatorios.
    \item El email no debe estar registrado previamente.
    \item Errores internos devuelven HTTP 500.
\end{itemize}

El login muestra errores genéricos ``Credenciales inválidas'' sin revelar
si el usuario existe o la contraseña es incorrecta.

\section{Tecnologías Utilizadas}

\begin{table}[h]
\centering
\caption{Stack tecnológico del proyecto}
\begin{tabularx}{\textwidth}{|l|X|}
\hline
\textbf{Tecnología} & \textbf{Versión} \\
\hline
Next.js & 15.5.18 \\
TypeScript & 5.8 \\
TailwindCSS & 4 \\
pnpm & 11.5 \\
Turborepo & 2.9 \\
PostgreSQL & 16 \\
Drizzle ORM & 0.41 \\
NextAuth & 5.0.0-beta.31 \\
Docker Compose & 3.8 \\
Node.js & 22 \\
\hline
\end{tabularx}
\end{table}

\section{Próximos Pasos}

\begin{enumerate}[leftmargin=*]
    \item \textbf{Sprint 2:} Gestión de Clientes — CRUD completo con tabla,
          formularios y búsqueda.
    \item \textbf{Sprint 3:} Gestión de Pacientes — Registro de mascotas
          vinculadas a clientes.
    \item \textbf{Sprint 4:} Agenda de Citas — Calendario y programación.
    \item \textbf{Sprint 5-6:} Historial clínico y módulo de facturación.
\end{enumerate}

\end{document}
